\documentclass[10pt,a4paper]{amsart}
\usepackage[margin=19mm]{geometry}
\usepackage{amsmath,amssymb,mathtools}
\usepackage[scr=rsfs,cal=euler]{mathalfa}
\usepackage{xfrac}
\usepackage[unicode,hidelinks,bookmarks=false]{hyperref}
\newcommand{\ie}{i.e.\ }
\newcommand{\cf}{cf.\ }
\newcommand{\resp}{respectively\ }
\newcommand{\Subsection}{\subsection}
\providecommand{\nobreakdash}{\nobreak\hbox{-}}
\setlength{\emergencystretch}{2em}
\multlinegap0pt
\usepackage{amssymb}
\usepackage{bm}
\newtheorem{prop}{Proposition}
\newtheorem{lem}{Lemma}
\newcommand{\Var}{\operatorname{Var}}
\newcommand{\ex}{\mathbb E}
\newcommand{\prim}{\operatorname{prim}}
\newcommand{\re}{\mathbb R}
\title{Thermal Concentration and Poisson--Dirichlet Edge Statistics for Random-Lattice Gibbs Ensembles}
\author{Masahiro Kaminaga}
\date{Source: arXiv:2607.00311v2 (CC BY 4.0)}
\begin{document}
\maketitle

\noindent\emph{Source and licence.} Thermal Concentration and Poisson--Dirichlet Edge Statistics for Random-Lattice Gibbs Ensembles. Version arXiv:2607.00311v2 (CC BY 4.0), distributed by arXiv under the Creative Commons Attribution 4.0 International licence, http://creativecommons.org/licenses/by/4.0/, as recorded in the arXiv licence metadata for this exact version and shown on the abstract page https://arxiv.org/abs/2607.00311v2. Full licence notice: \texttt{licenses/arxiv-2607.00311-CC-BY-4.0.txt}.\par\smallskip

\noindent\textit{Editorial scope.} Counted unit: the article's prop prop:local-critical-thermal-mass (source0.tex lines 1146--1170). The article's complete original proof of it is delivered with it (source0.tex lines 1172--1352, one \texttt{proof} environment of 181 lines). No mathematics is written, repaired or re-derived: the statement and the proof are the article's own lines, in the article's own notation, and no retained line is substituted or re-flowed.  Retained uncounted as the source-local context the proof needs: lines 764--784; lines 945--965.  Nothing was omitted from any retained span, so the package carries no omission record.  No retained line contains a citation, so the wrapper carries no bibliography and no bibliography entry.  No retained line contains a figure environment, an image-inclusion command or drawing code, so no image is delivered and the package declares no asset dependency.  Editorial formatting only, in addition: an AMS \texttt{amsart} preamble that restates the article's own declarations and macros used by the retained lines, namely the environments \texttt{prop}, \texttt{lem} on separate counters, together with the standard packages those lines need.  The retained environments print in this document as Proposition 1, Lemma 1 in order of appearance, whereas the article prints the same items with the numbering of its own full document.  The complete native source of the article as distributed in the arXiv e-print for this exact version is bundled unchanged as \texttt{sources/204163/source0.tex}.
\begin{prop}
\label{prop:weighted}
Assume $n\ge3$.
Let $f:\re^n\to[0,\infty)$ be an even Borel function such that $f$ and $f^2$ are integrable.
For $L\in\mathbb X_n$, define
$$
S_f(L)=\sum_{\bm u\in L_{\prim}^\sharp} f(\bm u),
$$
where $f(\bm u)$ means $f(\bm v)$ for either representative of $\bm u$.
Then, with respect to the random choice of $L$ according to $\mu_n$, the
mean and variance of $S_f(L)$ are given by 
$$
\ex_{\mu_n}S_f(L)=
\frac{1}{2\zeta(n)}\int_{\re^n}f(\bm x)\,d\bm x
$$
and
$$
\Var_{\mu_n}(S_f(L))=
\frac{1}{2\zeta(n)}\int_{\re^n}f(\bm x)^2\,d\bm x.
$$
\end{prop}

\begin{lem}
\label{lem:laplace}
Let $0<c<1$ and put
$$
y_c=\frac{1}{2}\log\frac{1}{c}
$$
and
$$
F(c)=\Phi_c(y_c)=
\frac{1}{2}\log\frac{1}{c}-\frac{1-c}{2}.
$$
Then $F(c)>0$.
Moreover, if
$$
\Psi_c(y)=y-c(e^{2y}-1),
$$
then
$$
\sup_{y\ge0}\Psi_c(y)<2F(c).
$$
\end{lem}

\begin{prop}
\label{prop:local-critical-thermal-mass}
Let $L_n$ be distributed according to $\mu_n$ on $\mathbb X_n$.
Assume $0<c<1$ and put
$$
y_c=\frac{1}{2}\log\frac{1}{c}.
$$
For $t\in\re$, define
$$
H(t)=\frac{1}{\sqrt{\pi}}\int_{-\infty}^{t}e^{-x^2}dx.
$$
This is the distribution function of the centered Gaussian law with density
$\pi^{-1/2}e^{-x^2}$, or equivalently of a normal random variable with
variance $1/2$.
Then
$$
\frac{
B_{n,c}((-\infty,y_c+t/\sqrt{n}],L_n)
}{
B_{n,c}(L_n)
}
\longrightarrow H(t)
$$
in probability as $n\to\infty$.
\end{prop}

\begin{proof}
We first consider the interval
$$
A_{n,t}=[0,y_c+t/\sqrt{n}].
$$
For all sufficiently large $n$ this interval is nonempty. Recall that
$$
J_{n,c}(A)=
n\int_A \exp\{n\Phi_c(y)\}dy
$$
and
$$
K_{n,c}(A)=
n\int_A \exp\{n\Psi_c(y)\}dy,
$$
where
$$
\Phi_c(y)=y-\frac{c}{2}(e^{2y}-1)
$$
and
$$
\Psi_c(y)=y-c(e^{2y}-1).
$$
By Proposition \ref{prop:weighted},
$$
\ex_{\mu_n}B_{n,c}(A,L)=\frac{J_{n,c}(A)}{2\zeta(n)}
$$
and
$$
\Var_{\mu_n}(B_{n,c}(A,L))
=
\frac{K_{n,c}(A)}{2\zeta(n)}.
$$

We compute the first moment near the maximum of $\Phi_c$.
The function $\Phi_c$ has its unique maximum on $[0,\infty)$ at $y_c$,
and
$$
\Phi_c(y_c)=F(c),\qquad \Phi_c''(y_c)=-2.
$$
Choose $\delta>0$ so small that $y_c-\delta>0$.
Since $y_c$ is the unique maximum of $\Phi_c$ on $[0,\infty)$, there is
$\eta>0$ such that
$$
\Phi_c(y)\le F(c)-\eta
$$
for all $y\in[0,\infty)$ with $|y-y_c|\ge\delta$.
Thus the contribution of this region to $J_{n,c}([0,\infty))$ and to
$J_{n,c}(A_{n,t})$ is exponentially smaller than $e^{nF(c)}$.

On the interval $|y-y_c|<\delta$, put
$$
y=y_c+\frac{x}{\sqrt{n}}.
$$
Then
$$
J_{n,c}(A_{n,t})
=
\sqrt{n}e^{nF(c)}
\int_{-\sqrt{n}y_c}^{t}
\exp\{n(\Phi_c(y_c+x/\sqrt{n})-F(c))\}dx
$$
up to an exponentially smaller error from the part outside
$(y_c-\delta,y_c+\delta)$.
Taylor's formula gives, locally uniformly in $x$,
$$
n(\Phi_c(y_c+x/\sqrt{n})-F(c))\longrightarrow -x^2 .
$$
Moreover, after decreasing $\delta$ if necessary, there is a constant
$\kappa>0$ such that
$$
\Phi_c(y_c+z)\le F(c)-\kappa z^2
$$
for $|z|\le\delta$.
Hence, in the variable $x=\sqrt n z$, the local integrand is bounded by
$e^{-\kappa x^2}$.
Therefore dominated convergence, together with the exponentially small
contribution outside the fixed neighborhood of $y_c$, gives
$$
J_{n,c}(A_{n,t})
=
\sqrt{n}e^{nF(c)}
\left\{
\int_{-\infty}^{t}e^{-x^2}dx+o(1)
\right\}.
$$
The same argument with the moving endpoint removed gives
$$
J_{n,c}([0,\infty))
=
\sqrt{n}e^{nF(c)}
\left\{
\int_{-\infty}^{\infty}e^{-x^2}dx+o(1)
\right\}
=
\sqrt{\pi n}e^{nF(c)}{1+o(1)}.
$$
Consequently,
$$
\frac{J_{n,c}(A_{n,t})}{J_{n,c}([0,\infty))}
\longrightarrow
\frac{1}{\sqrt{\pi}}\int_{-\infty}^{t}e^{-x^2}dx
=H(t).
$$
%
We next show that the corresponding random weighted sums are concentrated
around their means. By the elementary Laplace upper bound,
$$
\limsup_{n\to\infty}
\frac{1}{n}\log K_{n,c}([0,\infty))
\le
\sup_{y\ge0}\Psi_c(y).
$$
By Lemma \ref{lem:laplace}, the right hand side is strictly smaller than
$2F(c)$. Choose $G$ such that
$$
\sup_{y\ge0}\Psi_c(y)<G<2F(c).
$$
Then, for all sufficiently large $n$,
$$
K_{n,c}([0,\infty))\le e^{nG}.
$$
Since $J_{n,c}(A_{n,t})$ is of order $\sqrt n e^{nF(c)}$, Chebyshev's
inequality gives
$$
\frac{
B_{n,c}(A_{n,t},L_n)
}{
\ex_{\mu_n}B_{n,c}(A_{n,t},L)
}
\longrightarrow 1
$$
in probability. The same argument gives
$$
\frac{
B_{n,c}([0,\infty),L_n)
}{
\ex_{\mu_n}B_{n,c}([0,\infty),L)
}
\longrightarrow 1
$$
in probability.

It remains to remove the negative part of the $Y_n$--scale. For $y<0$,
the function $\Phi_c(y)$ has supremum $0$, whereas $F(c)>0$. Hence
$$
\limsup_{n\to\infty}
\frac{1}{n}\log J_{n,c}((-\infty,0))\le 0.
$$
By Markov's inequality,
$$
\frac{
B_{n,c}((-\infty,0),L_n)
}{
J_{n,c}([0,\infty))
}
\longrightarrow 0
$$
in probability. On the other hand, the preceding concentration estimate
implies that $B_{n,c}([0,\infty),L_n)$ is of order
$J_{n,c}([0,\infty))$ in probability. Therefore
$$
\frac{
B_{n,c}((-\infty,0),L_n)
}{
B_{n,c}([0,\infty),L_n)
}
\longrightarrow 0
$$
in probability. Thus the negative part of the $Y_n$--scale does not affect
the limiting ratio. Combining this fact with the Laplace ratio above gives
$$
\frac{
B_{n,c}((-\infty,y_c+t/\sqrt{n}],L_n)
}{
B_{n,c}(L_n)
}
\longrightarrow H(t)
$$
in probability. This proves the proposition.
\end{proof}

\end{document}
